\chapter{Regressão}

\section{Modelo Linear Geral}

\begin{definition}[Modelo Linear Geral]
	Um modelo linear geral amostral é dado por
	\begin{equation}
		Y_i = \beta_0 + \beta_1 x_{1i} \beta_1 x_{2i} + \cdots + \beta_p x_{pi} + \epsilon_i, \quad i=1,\ldots n, \quad n>p+1 
	\end{equation}
	\begin{outline}
		\1 $Y_i$: variável resposta, dependente ou predita (v.a.observável)
		\1 $x_{ji}$ variáveis regressoras, explicadoras, independentes ou preditoras (não aleatórias) ($j=1,\ldots,p$)
		\1 $\beta_j$ são os parâmetros ou coeficientes de regressão
			\2 Interpretação de $\beta_0$: 
				\3 $\E(Y)$
				\3 
		\1 $\epsilon_i$ sã v.a. não observáveis em que:
			\2 $\E(\epsilon_i)=0$
			\2 $\Var(\epsilon_i)=\sigma^2$
			\2 $\Cov(\epsilon_i, \epsilon_j)=\sigma_{ij}=0$, $i,j=1,\ldots,n$
	\end{outline}
	De forma matricial,
	\begin{align}
		\v{Y} &= \m{X} \v{\beta} + \v{\epsilon} \\
		\begin{bmatrix}
			Y_1 \\ Y_2 \\ \vdots \\ Y_n
		\end{bmatrix} &= 
		\begin{bmatrix}
			1 & x_{11} & x_{21} & \cdots & x_{p1} \\
			1 & x_{12} & x_{22} & \cdots & x_{p2} \\
			1 & \vdots & \vdots & & \vdots \\
			1 & x_{1n} & x_{2n} & \cdots & x_{pn} \\
		\end{bmatrix}
		\begin{bmatrix}
			\beta_1 \\ \beta_2 \\ \vdots \\ \beta_p
		\end{bmatrix} + 
		\begin{bmatrix}
			\epsilon_1 \\ \epsilon_2 \\ \vdots \\ \epsilon_n
		\end{bmatrix} 
	\end{align}
	\begin{outline}
		\1 $X$: matriz de planejamento
		\1 $\Var(\v{\epsilon}) = \m{\Sigma}$
	\end{outline}
\end{definition}

\begin{definition}[Definições complementares]
	\leavevmode
	\begin{outline}[enumerate]
		\1 Valores preditos: \( \hat{\v{Y}} = X\hat{\v{\beta}} \)
		\1 Resíduos: $\hat{\v{\epsilon}} = \v{Y} - \hat{\v{Y}}$
		\1 \( \m{H} = \m{X}(\m{X}\T\m{X})\inv \m{X}\T \)
	\end{outline}
\end{definition}

\section{Caso 1: \texorpdfstring{$\m{\Sigma} = \sigma^2 \v{I}_n$}{Sigma = sigma^2 I_n} e $\m{X}$ de posto completo}

Considere $\m{X}_{n \times (p+1)}$ de posto completo, ou seja $\pos(\m{X})=p+1$ e $\Var(\v{\epsilon})=\m{\Sigma}=\sigma^2\v{I}_n$

\begin{theorem}[Fatos úteis sobre $\m{H}$] Seja $\m{H} = \m{X}(\m{X}\T\m{X})\inv \m{X}\T$.
	\leavevmode
	\begin{outline}[enumerate]
		\1 $\m{H}$ é simétrica e idempotente
		\1 $\m{I}-\m{H}$ é simétrica e idempotente
		\1 $\m{H}\m{X}=\m{X}$
		\1 $\m{X}\T(\v{Y}-\hat{\v{Y}}) = \m{X}\T(\v{Y}-\m{H}\v{Y}) = \v{0}$
	\end{outline}
\end{theorem}

\begin{theorem}[Estimador de mínimos quadrados para $\v{\beta}$]
	\begin{equation}
		\hat{\v{\beta}}_{MQ} = (\m{X}\m{X}\T)\inv\m{X}\T \v{Y}
	\end{equation}
\end{theorem}

\begin{theorem}[Propriedades do estimador de mínimos quadrados]
	\leavevmode
	\begin{outline}[enumerate]
		\1 \( \E(\hat{\v{\beta}}_{MQ}) = \v{\beta} \)
		\1 \( \Var(\hat{\v{\beta}}_{MQ}) = \sigma^2(\m{X}\T\m{X})\inv \)
		\1 Se $\m{X}$ é de posto completo, $\v{a}\T\hat{\v\beta}$ é o melhor estimador linear não viesado (BLUE) de $\v{a}\T\v{\beta}$ para todo vetor $\v a$.
	\end{outline}
\end{theorem}

\begin{theorem}[Quadrado médio residual]
	O quadrado médio residual QMRes é um estimador não viesado para $\sigma^2$:
	\begin{equation}
		\hat{\sigma}^2 = \qmres \coloneq \frac{\sqres}{n-(p+1)} = \frac{{(\v{Y}-\hat{\v{Y}})\T(\v{Y}-\hat{\v{Y}})}}{n-(p+1)}
	\end{equation}
\end{theorem}

\begin{theorem}[Estimadores de máxima verossimilhança]
	Assumindo que $\v{\epsilon} \sim \normdist[n]{\v{0}, \sigma^2\v{I}_n}$, os estimadores de máxima verossimilhança para $\v{\beta}$ e $\sigma^2$ são:
	\begin{align}
		\hat{\v{\beta}}_{MV} &= \hat{\v{\beta}}_{MQ} = (\m{X}\m{X}\T)\inv\m{X}\T \v{Y} \\
		\hat{\sigma^2}_{MV} &= \frac{\sqres}{n} = \frac{(\v{Y}-\hat{\v{Y}})\T(\v{Y}-\hat{\v{Y}})}{n}
	\end{align}
\end{theorem}


\begin{theorem}[Propriedades dos estimadores de máxima verossimilhança]
	\leavevmode
	\begin{outline}[enumerate]
		\1 \( \hat{\v{\beta}}_{MV} \sim \normdist[p]{\v{\beta}, \sigma^2(\m{X}\T\m{X})\inv} \)
		\1 Suponha que queremos estimar $t(\v{\beta}, \sigma^2)$. Seja $q(\hat{\v{\beta}}, \hat{\sigma}^2)$ um ENV. Então, $q(\hat{\v{\beta}}, \hat{\sigma}^2)$ é ENVVUM. Consequências:
			\2 $\hat{\v{\beta}}$ é ENVVUM de $\v{\beta}$
			\2 $\hat{\sigma}^2$ é ENVVUM de $\sigma^2$
	\end{outline}
\end{theorem}

\section{Testes de hipóteses}

\subsection{Teste ANOVA}
 
 Considere o modelo linear geral do Caso 1 e seja
\begin{equation}
	\v{\beta} = \begin{bmatrix}
		\beta_ 0 \\ 
		\v{\beta}_1
	\end{bmatrix} = [\beta_0, \v{\beta}_1\T]\T
\end{equation} 
 Queremos avaliar as hipóteses
 
 \begin{equation}
 	\left\{ \begin{array}{ll} 
 		H_0: & \v{\beta}_1 = \v{0} \\
 		H_1: & \text{pelo menos uma desigualdade}.
 	\end{array} \right.
 \end{equation}
       
 
\subsubsection{Modelo com intercepto}

Para o modelo com intercepto, a tabela ANOVA fica

\begin{table}[h]
	\centering
	\begin{tblr}{
			colspec = {|c|c|c|c|c|},
			cells   = {valign=m, halign=c},
			rowsep  = 3pt,
			hlines, vlines,
		}
		F.V. & g.l. & SQ & QM & F \\
		Regressão & $p$ & $\sqr$ & $\displaystyle\frac{\sqr}{p}$ &
		\SetCell[r=2]{} $\displaystyle \frac{\qmr}{\qmres}$ \\
		Resíduos & $n-(p+1)$ & $\sqres$ & $\displaystyle\frac{\sqres}{n-(p+1)}$ & \\
		Total & $n-1$ & $\sqt$ & -- & -- \\
	\end{tblr}
	\caption{ANOVA para o modelo com intercepto}
\end{table}

Então, 

\begin{equation}\displaystyle
	F = \frac{\frac{\sqr}{p}}{\frac{\sqres}{n-(p+1)}} \undernull \fdist{p, n-(p+1)}
\end{equation}

Rejeitamos $H_0$ se $f > f_c$, em que $\P(f>f_c)=\alpha$.

\begin{align}
	\sqt &= \sqres + \sqr  \\
	\sum_{i=1}^n (Y_i - \bar{Y})^2 &= 
	\sum_{i=1}^n (Y_i - \hat{Y}_i)^2 +
	\sum_{i=1}^{n}(\hat{Y}_i - \bar{Y})^2 + 
\end{align}

\begin{align}
	\sqt &= \sum_{i=1}^n (Y_i - \bar{Y})^2 = \sum_{i=1}^{n} Y_i^2 -n\bar{Y}^2 \\
	&= \v{Y}\T\left( \m{I} - \frac{\v{1}\v{1}\T}{n} \right) \v{Y}
\end{align}

\begin{align}
	\sqres &= \sum_{i=1}^n (Y_i - \hat{Y}_i)^2 \\
	&= (\v{Y}-\hat{\v{Y}})\T (\v{Y}-\hat{\v{Y}}) \\
	&= \v{Y}\T (\m{I}-\m{H}) \v{Y} \\
	&= \v{Y}\T \v{Y} - \hat{\v{\beta}}\T \m{X}\T \v{Y}
\end{align}

\begin{align}
	\sqr &= \sum_{i=1}^{n}(\hat{Y}_i - \bar{Y})^2 \\
	&=  \v{Y}\T\left( \m{H} - \frac{\v{1}\v{1}\T}{n} \right) \v{Y}
\end{align}

\begin{theorem}[Distribuições no modelo com intercepto]
	Assumindo que $\m{X}$ é de posto completo, com $\pos(\m{X}) = p+1$.
	\leavevmode
	\begin{outline}[enumerate]
		\1 \(\displaystyle \frac{\sqres}{\sigma^2} \sim \chidist{n-(p+1)}  \)
		\1 \(\displaystyle \frac{\sqr}{\sigma^2} \sim \chidist{p, \lambda} \)
			\2 \(\displaystyle \lambda = \frac{1}{2\sigma^2}\v{\beta}_1\T \m{X}_1\T \left( \m{I} - \frac{\v{1}\v{1}\T}{n} \right) \), onde $\v{\beta} = [\beta_0, \ \v{\beta}_1\T]\T$.
			\2 Sob $H_0$: $\v{\beta}_1=0 \implies \lambda=0$.
		\1 \( \displaystyle \frac{\sqt}{\sigma^2} \sim \chidist{n, \lambda}  \), com $\lambda$ igual ao item anterior.
		\1 $\sqres$ e $\sqr$ são independentes.
	\end{outline}
\end{theorem}


\subsubsection{Modelo sem intercepto}

No modelo sem intercepto, a decomposição usual de $\sqt$ não é válida. Desta forma, é necessário trabalhar com novas quantidades não corrigidas pela média.

Note agora que $\pos(\m X) = p$.

\begin{table}[h]
	\centering
	\begin{tblr}{
			colspec = {|c|c|c|c|c|},
			cells   = {valign=m, halign=c},
			rowsep  = 3pt,
			hlines, vlines,
		}
		F.V. & g.l. & SQ & QM & F \\
		Regressão n.c. & $p$ & $\sqr_{nc}$ & $\displaystyle\frac{\sqr_{nc}}{p}$ &
		\SetCell[r=2]{} $\displaystyle \frac{\qmr}{\qmres}$ \\
		Resíduos & $n-p$ & $\sqres$ & $\displaystyle\frac{\sqres}{n-p}$ & \\
		Total n.c. & $n$ & $\sqt_{nc}$ & -- & -- \\
	\end{tblr}
	\caption{ANOVA para o modelo sem intercepto}
\end{table}


\begin{equation}\displaystyle
	F = \frac{\frac{\sqr_{nc}}{p}}{\frac{\sqres}{n-p}} \undernull \fdist{p, n-p}
\end{equation}

Rejeitamos $H_0$ se $f > f_c$, em que $\P(f>f_c)=\alpha$.


\begin{align}
	\sqt_{nc} &= \sqres + \sqr_{nc}  \\
	\sum_{i=1}^n Y_i^2 &= 
	\sum_{i=1}^n (Y_i-\hat{Y}_i)^2 +
	\sum_{i=1}^n \hat{Y}_i^2 + 
\end{align}

\begin{align}
	\sqt_{nc} &= \sum_{i=1}^n Y_i^2 \\
	&= \v{Y}\T \v{Y}
\end{align}

\begin{align}
	\sqres &= \sum_{i=1}^n (Y_i - \hat{Y}_i)^2 \\
	&= (\v{Y}-\hat{\v{Y}})\T (\v{Y}-\hat{\v{Y}}) \\
	&= \v{Y}\T (\m{I}-\m{H}) \v{Y} \\
	&= \v{Y}\T \v{Y} - \hat{\v{\beta}}\T \m{X}\T \v{Y}
\end{align}

\begin{align}
	\sqr_{nc} &= \sum_{i=1}^n \hat{Y}_i^2 \\
	&= \v{Y}\T \m{H} \v{Y} = \hat{\v{\beta}}\T\m{X}\T\v{Y}
\end{align}

\begin{theorem}[Distribuições no modelo sem intercepto]
	Assumindo que $\m{X}$ é de posto completo com $\pos(\m{X}) = p$.
	\leavevmode
	\begin{outline}[enumerate]
		\1 \(\displaystyle \frac{\sqres}{\sigma^2} \sim \chidist{n-p}  \)
		\1 \(\displaystyle \frac{\sqr}{\sigma^2} \sim \chidist{p, \lambda} \)
		\2 \(\displaystyle \lambda = \frac{1}{2\sigma^2}\v{\beta}\T \m{X}\T \m{X} \v{\beta}\)
		\2 Sob $H_0$: $\v\beta=0 \implies \lambda=0$.
		\1 \( \displaystyle \frac{\sqt}{\sigma^2} \sim \chidist{n, \lambda}  \), com $\lambda$ igual ao item anterior.
		\1 $\sqres$ e $\sqr$ são independentes.
	\end{outline}
\end{theorem}



\subsection{Teste da Razão de Verossimilhanças Generalizado}

Queremos avaliar a hipótese linear geral:

\begin{equation}
	\left\{ \begin{array}{ll} 
		H_0: & \m{C} \v\beta = \v{m} \\
		H_1: & \m{C} \v\beta \neq \v{m} 
	\end{array} \right.
\end{equation}
%
onde $\m{C}$ é uma matriz dedimensão $q \times (p+1)$ com $\pos(\m{C})=q<p+1$ e $\v{m}$ um vetor de dimensão $q \times 1$.

\begin{theorem}[Estimadores de MV restritos a $H_0$] 
	Sob $H_0$, os estimadores de máxima verossimilhança restritos são dados por
	\begin{align}
		\hat{\v{\beta}}_0 &= \hat{\v{\beta}} - (\m{X}\T\m{X})\inv \m{C}\T [\m{C} (\m{X}\T\m{X})\inv\m{C}\T]\inv(\m{C}\hat{\v{\beta}}-\v{m})  \\
		\hat{\sigma}^2_0 &= \frac{(\v{Y}-\m{X}\hat{\v{\beta}}_0)\T(\v{Y}-\m{X}\hat{\v{\beta}}_0)}{n}
	\end{align}
\end{theorem}

\begin{theorem}[TRV]
	Seja
	
	\begin{equation}
		W = \frac{(\m{C}\hat{\v{\beta}}-\v{m})\T [\m{C}(\m{X}\T \m{X})\inv \m{C}\T]\inv (\m{C}\hat{\v{\beta}}-\v{m})}{q\hat{\sigma}^2}
	\end{equation}
		%
	em que
	\begin{equation}
		\hat{\sigma}^2 = \frac{(\v{Y}-\m{X}\hat{\v{\beta}})\T(\v{Y}-\m{X}\hat{\v{\beta}})}{n-(p+1)}.
	\end{equation}
	Então,
	\begin{equation}
		W \undernull \fdist{q,n-(p+1)}
	\end{equation}
	%
	e rejeitamos $H_0$ ao nível de significância $\alpha$ se $w > w_0$, com $w_0$ tal que $\P(W > w_0)=\alpha$.
\end{theorem}

\begin{theorem}[Fórmulas alternativas para $W$]
	\leavevmode
	\begin{align}
		W &= \left(\frac{\hat{\sigma}^2_0-\hat{\sigma}^2}{\hat{\sigma}^2} \right) \frac{n-(p+1)}{q} \\
		&= \left( \frac{\sqres_0-\sqres}{\sqres} \right) \frac{n-(p+1)}{q}
	\end{align}
\end{theorem}


\begin{theorem}[Poder do teste]
	Cada $\v{\beta}$ associado a uma hipótese $H_1$ do teste $\Upsilon$ corresponde a um $\lambda$ diferente. Sob $H_1$,
	
	\begin{equation}
		W = \frac{\sqr/p}{\sqres/(n-p-1)} \underalt \fdist{p, n-p-1, \lambda}
	\end{equation}
	
	\begin{align}
		\pi_{\Upsilon} = \pi(\lambda) &= \P[\text{rejeitar } H_0 \mid H_1] \\
		&= \int_{w_0}^{\infty} f_{q,n-p-1,\lambda}  dw \\
		&= 1-F_{q,n-p-1,\lambda}(w_0) \\
		&= 1-\texttt{pf}(w_0, q, n-p-1, \texttt{nc}=2\lambda)
	\end{align}
	
\end{theorem}

\section{Regressão linear simples}


\begin{equation}
	Y_i = \beta_0 + \beta_1 x_i + \epsilon_i, \quad i=1,\ldots,n
\end{equation}
%
em que $\E(\epsilon_i)=0$ e $\Var(\epsilon_i)=\sigma^2$ e $\epsilon_i, \epsilon_j$ são não correlacionados.

Em notação matricial:

\begin{outline}
	\1 $\v{Y} = (Y_1, \ldots, Y_n)\T$ 
	\1 \( \m{X} = \begin{bmatrix}
		1 & x_1 \\ 1 & x_2 \\ \vdots & \vdots \\ 1 & x_n
	\end{bmatrix} \)
	\1 \(\m{X}\m{X}\T = \begin{bmatrix}
		n & \sum_{i=1}^n x_1  \\ \sum_{i=1}^n x_i & \sum_{i=1}^n x_i^2
	\end{bmatrix} \) 
	\1 \( \displaystyle \ (\m{X}\m{X}\T)\inv = \frac{1}{n\sum_{i=1}^n x_i^2 - (\sum_{i=1}^n x_i)^2} \begin{bmatrix}
		\sum_{i=1}^n x_i^2 & -\sum_{i=1}^n x_i  \\ -\sum_{i=1}^n x_i & n
	\end{bmatrix} \)
	\1 \( \m{X}\T\m{Y} = \begin{bmatrix}
	\sum_{i=1}^n Y_i  \\ \sum_{i=1}^n x_i Y_i
	\end{bmatrix} \)
\end{outline}


Definindo
\begin{equation}
	S_{xx} = \sum_{i=1}^n (x_i-\bar{x})^2, \quad S_{xY} = \sum_{i=1}^n (x_i-\bar{x})(Y_i-\bar{Y})
\end{equation}
%
temos que

\begin{equation}
	\hat{\v{\beta}} = \begin{bmatrix} 
		\hat{\beta}_0 \\ \hat{\beta}_1
	\end{bmatrix}
	= (\m{X}\T \m{X})\inv\m{X}\T\v{Y} = 
	\begin{bmatrix}
		\bar{Y} - \hat{\beta_1}\bar{x} \\ \displaystyle \frac{S_{xY}}{S_{xx}}
	\end{bmatrix}
\end{equation}

Forma alternativa:

\begin{align}
	\hat{\beta}_1 = \frac{\sum_{i=1}^n x_i Y_i - n \bar{x}\bar{Y}}{\sum_{i=1}^{n}x_i^2-n \bar{x}^2} 
\end{align}